%%%%%%%%%%  scan idx 189  |  folio 176  %%%%%%%%%%
la
sous-catégorie triangulée $\mathrm D^b(A)_{\coh}$ (resp. $\mathrm D(A)_{\parf}$) formée des complexes à cohomologie bornée et pseudo-cohérents (resp. parfaits) relativement à la sous-catégorie $\Lib(A)$ (resp. $\Loclib(A)$) des modules libres (resp. projectifs) de type fini. On a des diagrammes commutatifs d'équivalences de catégories
\[
\begin{tikzcd}
\mathrm K^{-,b}(\Lib(S)) \arrow[r] \arrow[d,"\Gamma_S"] &
\mathrm K^{-,b}(\Loclib(S)) \arrow[r] \arrow[d,"\Gamma_S"] &
\mathrm D^b(S)_{\coh} \arrow[d,"\mathbb R\Gamma_S"] \\
\mathrm K^{-,b}(\Lib(A)) \arrow[r] &
\mathrm K^{-,b}(\Loclib(A)) \arrow[r] &
\mathrm D^b(A)_{\coh},
\end{tikzcd}
\]
\[
\begin{tikzcd}
\mathrm K^b(\Loclib(S)) \arrow[r] \arrow[d,"\Gamma_S"] &
\mathrm D(S)_{\parf} \arrow[d,"\mathbb R\Gamma_S"] \\
\mathrm K^b(\Loclib(A)) \arrow[r] &
\mathrm D(A)_{\parf},
\end{tikzcd}
\]
(où, dans le diagramme du haut, on peut remplacer $b$ par $-$ quand $S$ est de dimension cohomologique finie au sens de (3.7)).

b) Soit $S$ un schéma. Il résulte de a) qu'un complexe de $\mathcal O_S$-Modules pseudo-cohérent et à cohomologique bornée est localement isomorphe (dans $\mathrm D(S)$) à un complexe strictement pseudo-cohérent, ce qui n'était nullement évident a priori (cf. (I 2.3.1)).

\subsection*{2.3. Faisceaux mous}

\textbf{2.3.1.} Rappelons quelques définitions et résultats de [T.F]. Soit $X$ un espace paracompact. On dit qu'un faisceau d'ensemble $F$ sur $X$ est \emph{mou} si, pour tout fermé $Y$ de $X$, toute section de $F$ au-dessus de $Y$ se prolonge en une section de $F$ au-dessus de $X$. Les faisceaux abéliens mous sont acycliques, i.e. si $F$ est mou, on a $H^i(X,F)=0$ pour $i>0$ (T.F. II 4.4.3).

Soit $A$ un faisceau d'anneaux sur $X$. Si $F$ est un $A$-Module mou, $F$ est engendré par ses sections (cf. 2.1.1) : il est clair en effet que l'homomorphisme $A_X^{(I)}\to F$ défini par $I=H^0(X,F)$ est un épimorphisme. 
